% GENERATED FILE. Edit canonical lessons, then run npm run generate:books.
% canonical-source-hash: fnv1a64:5e132b216d5f1004
% canonical-lessons: ES-C613-requerir, ES-C613-requisito, ES-C613-incertidumbre, ES-C613-probable, ES-C613-probabilidad

\chapter{Lo que hace falta}
\label{ch:es-lo-que-hace-falta}
\hlchaptermodality{613}

\begin{chapteropening}
\textbf{By the end of this chapter:} \emph{I can say what something requires and how likely it is.}

Complete the last lesson of chapter 613: I can say what something requires and how likely it is.
\end{chapteropening}

\section[requerir]{requerir — to require, to call for}

\label{lesson:ES-C613-requerir}



\begin{quote}
Before the new one: say the Spanish for in that case, then the Spanish for in any case.
\end{quote}

\subsection*{You'll want to know: requerir}
\textbf{requerir} — \textquotedblleft{}to require, to call for\textquotedblright{}.

\begin{quote}
El puesto requiere experiencia.
\end{quote}

\begin{quote}
\emph{The post requires experience.}
\end{quote}

Like \emph{sentir}: \emph{requiere, requirió}. Common in job offers and rules.

\begin{cousinweb}
\emph{requerir} is inherited from Latin \emph{requirere} 'to look for again, to need', built from \emph{re} and \emph{quaerere} 'to seek', and it follows \emph{querer} in its forms (EtymDB).
\end{cousinweb}

\begin{grammarlens}[title={What you've built}]
One more word for telling, explaining and describing at B1.
\end{grammarlens}

\subsection*{Your turn}
\begin{itemize}
\raggedright
  \item \emph{Say it:} \emph{requerir}
  \item \emph{Say it:} \emph{requerir}, once more
  \item \emph{Say it:} a sentence of your own with \emph{requerir}
  \item \emph{Recall:} say the Spanish for in that case, then the Spanish for in any case, then say \emph{requerir} again
\end{itemize}

\begin{tcolorbox}[breakable,colback=teal!4,colframe=teal!35!black,title={Before you move on}]
What does requerir mean? (\textquotedblleft{}To require, to call for\textquotedblright{}.) Say it once more.
\end{tcolorbox}
\section[el requisito]{el requisito — the requirement}

\label{lesson:ES-C613-requisito}



\begin{quote}
Before the new one: say the Spanish for in any case, then the Spanish for to require.
\end{quote}

\subsection*{You'll want to know: el requisito}
\textbf{el requisito} — \textquotedblleft{}the requirement\textquotedblright{}.

\begin{quote}
El inglés es un requisito para el puesto.
\end{quote}

\begin{quote}
\emph{English is a requirement for the post.}
\end{quote}

\emph{Cumplir los requisitos} — to meet the requirements.

\begin{cousinweb}
\emph{requisito} is a learned borrowing from Latin \emph{requisitus} 'required', the participle of \emph{requirere}, built on \emph{quaerere} 'to seek' (EtymDB).
\end{cousinweb}

\begin{grammarlens}[title={What you've built}]
One more word for telling, explaining and describing at B1.
\end{grammarlens}

\subsection*{Your turn}
\begin{itemize}
\raggedright
  \item \emph{Say it:} \emph{el requisito}
  \item \emph{Say it:} \emph{el requisito}, once more
  \item \emph{Say it:} a sentence of your own with \emph{el requisito}
  \item \emph{Recall:} say the Spanish for in any case, then the Spanish for to require, then say \emph{el requisito} again
\end{itemize}

\begin{tcolorbox}[breakable,colback=teal!4,colframe=teal!35!black,title={Before you move on}]
What does el requisito mean? (\textquotedblleft{}The requirement\textquotedblright{}.) Say it once more.
\end{tcolorbox}
\section[la incertidumbre]{la incertidumbre — the uncertainty}

\label{lesson:ES-C613-incertidumbre}



\begin{quote}
Before the new one: say the Spanish for to require, then the Spanish for the requirement.
\end{quote}

\subsection*{You'll want to know: la incertidumbre}
\textbf{la incertidumbre} — \textquotedblleft{}the uncertainty\textquotedblright{}.

\begin{quote}
Hay mucha incertidumbre sobre el futuro.
\end{quote}

\begin{quote}
\emph{There is a lot of uncertainty about the future.}
\end{quote}

The opposite of \emph{la certeza}.

\begin{cousinweb}
\emph{incertidumbre} is usually traced to Late Latin \emph{incertitūdō} 'uncertainty', a learned word built on \emph{incertus} 'not sure', from \emph{certus} 'sure' (Lewis \& Short).
\end{cousinweb}

\begin{grammarlens}[title={What you've built}]
One more word for telling, explaining and describing at B1.
\end{grammarlens}

\subsection*{Your turn}
\begin{itemize}
\raggedright
  \item \emph{Say it:} \emph{la incertidumbre}
  \item \emph{Say it:} \emph{la incertidumbre}, once more
  \item \emph{Say it:} a sentence of your own with \emph{la incertidumbre}
  \item \emph{Recall:} say the Spanish for to require, then the Spanish for the requirement, then say \emph{la incertidumbre} again
\end{itemize}

\begin{tcolorbox}[breakable,colback=teal!4,colframe=teal!35!black,title={Before you move on}]
What does la incertidumbre mean? (\textquotedblleft{}The uncertainty\textquotedblright{}.) Say it once more.
\end{tcolorbox}
\section[probable]{probable — likely, probable}

\label{lesson:ES-C613-probable}



\begin{quote}
Before the new one: say the Spanish for the requirement, then the Spanish for the uncertainty.
\end{quote}

\subsection*{You'll want to know: probable}
\textbf{probable} — \textquotedblleft{}likely, probable\textquotedblright{}.

\begin{quote}
Es probable que llueva.
\end{quote}

\begin{quote}
\emph{It is likely to rain.}
\end{quote}

\emph{Es probable que} + subjunctive.

\begin{cousinweb}
\emph{probable} is a learned borrowing from Latin \emph{probabilis} 'credible, that can be proved', built on \emph{probare} 'to test' (EtymDB).
\end{cousinweb}

\begin{grammarlens}[title={What you've built}]
One more word for telling, explaining and describing at B1.
\end{grammarlens}

\subsection*{Your turn}
\begin{itemize}
\raggedright
  \item \emph{Say it:} \emph{probable}
  \item \emph{Say it:} \emph{probable}, once more
  \item \emph{Say it:} a sentence of your own with \emph{probable}
  \item \emph{Recall:} say the Spanish for the requirement, then the Spanish for the uncertainty, then say \emph{probable} again
\end{itemize}

\begin{tcolorbox}[breakable,colback=teal!4,colframe=teal!35!black,title={Before you move on}]
What does probable mean? (\textquotedblleft{}Likely, probable\textquotedblright{}.) Say it once more.
\end{tcolorbox}
\section[la probabilidad]{la probabilidad — the probability, the likelihood}

\label{lesson:ES-C613-probabilidad}



\begin{quote}
Before the new one: say the Spanish for the uncertainty, then the Spanish for likely.
\end{quote}

\subsection*{You'll want to know: la probabilidad}
\textbf{la probabilidad} — \textquotedblleft{}the probability, the likelihood\textquotedblright{}.

\begin{quote}
La probabilidad es baja.
\end{quote}

\begin{quote}
\emph{The likelihood is low.}
\end{quote}

\emph{Con toda probabilidad} — in all likelihood.

\begin{cousinweb}
\emph{probabilidad} is built on \emph{probable}, a learned borrowing from Latin \emph{probabilis}, from \emph{probare} 'to test' (Lewis \& Short).
\end{cousinweb}

\begin{grammarlens}[title={What you've built}]
One more word for telling, explaining and describing at B1.
\end{grammarlens}

\subsection*{Your turn}
\begin{itemize}
\raggedright
  \item \emph{Say it:} \emph{la probabilidad}
  \item \emph{Say it:} \emph{la probabilidad}, once more
  \item \emph{Say it:} a sentence of your own with \emph{la probabilidad}
  \item \emph{Recall:} say the Spanish for the uncertainty, then the Spanish for likely, then say \emph{la probabilidad} again
\end{itemize}

\begin{tcolorbox}[breakable,colback=teal!4,colframe=teal!35!black,title={Before you move on}]
What does la probabilidad mean? (\textquotedblleft{}The probability, the likelihood\textquotedblright{}.) Say it once more.
\end{tcolorbox}
